\section*{Bab \ref{ch:b3:exploration-seismology} --- Seismik Eksplorasi}

\begin{solution}{exo:b3:exploration-seismology:1}
$Z = \rho v$: kapur $2000\times2500 = \qty{5.00e6}{kg.m^{-2}.s^{-1}}$, serpih
$2300\times2300 = \qty{5.29e6}{kg.m^{-2}.s^{-1}}$, batu gamping $2600\times4500 =
\qty{11.70e6}{kg.m^{-2}.s^{-1}}$. Kapur--serpih: $R = (5.29 - 5.00)/(5.29 + 5.00)
= 0.29/10.29 = 0.0282$; serpih--batu gamping: $R = (11.70 - 5.29)/(11.70 + 5.29) =
6.41/16.99 = 0.377$. Energi terpantul adalah $R^2 = 0.142$: bagian atas
batu gamping memantulkan \qty{14.2}{\percent} energi datang, antarmuka
kapur--serpih kurang dari \qty{0.1}{\percent}.
\end{solution}

\begin{solution}{exo:b3:exploration-seismology:2}
$t_0 = 2\times800/2000 = \qty{0.800}{s}$; $t = \sqrt{0.800^2 + (1000/2000)^2} =
\sqrt{0.89} = \qty{0.9434}{s}$; moveout-nya adalah \qty{0.1434}{s}. Pendekatan
memberikan $1000^2/(2\times2000^2\times0.800) = \qty{0.1563}{s}$, terlalu
tinggi sebesar \qty{9}{\percent}: pada offset yang lebih besar daripada kedalaman, $x \ll vt_0$ tidak
lagi berlaku.
\end{solution}

\begin{solution}{exo:b3:exploration-seismology:3}
Titik tengah bergeser sebesar setengah jarak penerima: interval CMP adalah
$12.5/2 = \qty{6.25}{m}$. \omterm{def:b3:exploration-seismology:cmp}{Liputan} adalah $240\times12.5/(2\times25) = 60$. \omterm{def:b3:exploration-seismology:cmp}{Liputan}
sebesar 120 memerlukan tembakan setiap $240\times12.5/(2\times120) = \qty{12.5}{m}$.
\end{solution}

\begin{solution}{exo:b3:exploration-seismology:4}
$\lambda = v/f$: \qty{100}{m} di batu pasir, $\lambda/4 = \qty{25}{m}$;
\qty{150}{m} di batu gamping, $\lambda/4 = \qty{37.5}{m}$. Dengan
$\lambda = 2800/30 = \qty{93.3}{m}$ dan $z = \qty{1000}{m}$,
$r_F \approx \sqrt{\lambda z/2} = \sqrt{46\,667} = \qty{216}{m}$ (\qty{217}{m}
dengan suku $\lambda^2/16$).
\end{solution}

\begin{solution}{exo:b3:exploration-seismology:5}
\omterm{thm:b3:exploration-seismology:dix}{Persamaan Dix}: $v^2 = (2800^2\times0.700 - 2400^2\times0.500)/0.200 = (5\,488\,000 -
2\,880\,000)/0.200 = \qty{1.304e7}{m^2.s^{-2}}$, maka $v = \qty{3611}{m.s^{-1}}$.
Lapisan itu dilintasi dalam \qty{0.100}{s} satu arah: $h = 3611\times0.100 =
\qty{361}{m}$.
\end{solution}

\begin{solution}{exo:b3:exploration-seismology:6}
(a) Misalkan $d$ adalah jarak dari titik tengah M ke reflektor, sehingga
$t_0 = 2d/v$. Sinar terpantul dari sumber S ke penerima G sama panjang
dengan garis lurus dari S$'$, bayangan S di reflektor, ke G. Letakkan M pada
titik asal, garis sepanjang kemiringan; S$'$ berada pada jarak $2d_S$ dari S sepanjang
normal, dengan $d_S = d - (x/2)\sin\varphi$. Kuadrat S$'$G adalah
$x^2 + 4d_S^2 + 4d_Sx\sin\varphi = x^2 + 4d^2 - x^2\sin^2\varphi$, semua suku
dalam $dx$ saling meniadakan, sehingga $v^2t^2 = 4d^2 + x^2\cos^2\varphi$ dan
$t^2 = t_0^2 + x^2\cos^2\varphi/v^2$. (b) Kumpulan jejak mengikuti hiperbola dengan kecepatan
$V = v/\cos\varphi = 2000/\cos\ang{12} = \qty{2045}{m.s^{-1}}$: kemiringan membuat
\omterm{def:b3:exploration-seismology:velocities}{kecepatan stacking} terlalu tinggi, di sini sebesar \qty{2.2}{\percent}.
\end{solution}

\begin{solution}{exo:b3:exploration-seismology:7}
Dengan $\sqrt{60} = 7.75$. Menggandakan rasio sinyal terhadap derau memerlukan
empat kali lebih banyak jejak: \omterm{def:b3:exploration-seismology:cmp}{liputan} sebesar 240.
\end{solution}

\begin{solution}{exo:b3:exploration-seismology:8}
Kedalaman yang diplot $vt/2$ berubah sebesar $\qty{1.5}{km.s^{-1}}\times\qty{0.40}{s}
= \qty{0.60}{km}$ setiap kilometer: $\tan\theta_a = 0.60$, $\theta_a = \ang{31.0}$.
Menurut teorema zero-offset $\sin\theta = 0.60$: $\theta = \ang{36.9}$.
\end{solution}

\begin{solution}{exo:b3:exploration-seismology:9}
Setiap lonjakan menempatkan salinan $w$ yang berpusat padanya: $0.3\,w$ pada sampel 0--2 menghasilkan
$(-0.15, 0.3, -0.15)$, dan $-0.2\,w$ pada sampel 2--4 memberikan $(0.1, -0.2, 0.1)$.
Jumlahnya adalah $(0, -0.15, 0.3, -0.05, -0.2, 0.1, 0)$, ditulis dari satu sampel
sebelum deret sampai satu sampel sesudahnya. Refleksi memuncak pada sampel 1 dan 3, dengan
tanda sebenarnya; keduanya tumpang tindih pada sampel 2, di mana $-0.15 + 0.1 = -0.05$: sedikit
interferensi, tetapi dua puncak yang terpisah.
\end{solution}

\begin{solution}{exo:b3:exploration-seismology:10}
$\arcsin(1800/2900) = \ang{38.37}$ dan $\arcsin(1800/4400) = \ang{24.15}$, maka
$i_c = (38.37 + 24.15)/2 = \ang{31.26}$ dan $\theta = (38.37 - 24.15)/2 =
\ang{7.11}$; $v_2 = 1800/\sin\ang{31.26} = \qty{3469}{m.s^{-1}}$. Rata-rata
harmonik $2\times2900\times4400/(2900 + 4400) = \qty{3496}{m.s^{-1}}$ adalah
\qty{0.8}{\percent} lebih tinggi; dikalikan dengan $\cos\theta = 0.9923$, diperoleh
\qty{3469}{m.s^{-1}}, pendekatan kemiringan kecil yang lazim
$v_2 \approx \cos\theta\,2v_dv_u/(v_d + v_u)$.
\end{solution}

\begin{solution}{exo:b3:exploration-seismology:11}
(a) Kompositnya adalah $R[w(t) - w(t - \Delta t)] \approx R\,\Delta t\,w'(t)$ dengan
ekspansi Taylor orde pertama. (b) Amplitudonya sebanding dengan
$\Delta t$, sehingga lapisan yang lebih tipis daripada \omterm{def:b3:exploration-seismology:resolution}{ketebalan tuning} tetap memantulkan sinyal,
dengan bentuk $w'$, namun bagian atas dan dasarnya tidak dapat dipilih sebagai dua peristiwa: ia
terdeteksi, tetapi tidak dipisahkan. Jika $R$ dan wavelet diketahui, misalnya dari
sumur, amplitudonya memberi $\Delta t$ dan karenanya ketebalannya. (c) Amplitudonya
sebanding dengan ketebalan: rasionya adalah $5/10 = 1/2$.
\end{solution}

\begin{solution}{exo:b3:exploration-seismology:12}
(a) $t(x) = (2/v)\sqrt{z^2 + x^2}$: $t(0) = 2\times600/3000 = \qty{0.400}{s}$ dan
$t(300) = (2/3000)\sqrt{600^2 + 300^2} = \qty{0.4472}{s}$. (b) $t^2 = t_0^2 +
x^2/(v/2)^2$, bentuk NMO dengan kecepatan $v/2$. Refleksi dengan waktu
puncak yang sama pada penampang zero-offset tidak akan menunjukkan moveout sama sekali; dibandingkan
dengan hiperbola NMO yang memiliki $t_0$ sama dalam kumpulan jejak, moveout difraksi
$x^2/(2(v/2)^2t_0)$ empat kali lebih besar, karena waktu bergantung pada
jarak ke penghambur dua kali, ke bawah dan kembali. (c) Migrasi menjumlahkan
penampang sepanjang hiperbola ini dan menempatkan energi di puncaknya: difraksi
runtuh menjadi titik pada $x = 0$, $z = \qty{600}{m}$.
\end{solution}

\begin{solution}{pb:b3:exploration-seismology:1}
\textbf{1.} Dari B, $400/(486.1 - 425.2) = 400/0.0609 = \qty{6568}{m.s^{-1}}$;
dari T, $400/(486.1 - 372.4) = 400/0.1137 = \qty{3518}{m.s^{-1}}$. Gelombang kepala
dari B tampak cepat: B adalah ujung bawah searah kemiringan, dan tembakannya merambat ke atas kemiringan;
refraktor naik dari B ke T.

\textbf{2.} $\arcsin(2400/3518) = \ang{43.02}$, $\arcsin(2400/6568) =
\ang{21.43}$: $i_c = \ang{32.22}$, $\theta = \ang{10.79}$.

\textbf{3.} $v_2 = 2400/\sin\ang{32.22} = \qty{4501}{m.s^{-1}}$: berada dalam
kisaran batu gamping, yaitu puncak batu gamping Jura.

\textbf{4.} Waktu intersep adalah $t - x/v_{\text{app}}$ pada \qty{1200}{m}:
$t_{i,B} = 0.4861 - 1200/6568 = \qty{0.3034}{s}$ dan $t_{i,T} = 0.4861 -
1200/3518 = \qty{0.1450}{s}$. Dengan $\cos i_c = 0.8460$, $h = v_1t_i/(2\cos
i_c)$: $h_W = 2400\times0.3034/1.692 = \qty{430.4}{m}$ dan $h_E =
\qty{205.7}{m}$.

\textbf{5.} Dibagi dengan $\cos\theta = 0.9823$: \qty{438.1}{m} di bawah B dan
\qty{209.4}{m} di bawah T. Maka $(438.1 - 209.4)/1200 = 0.1906 =
\tan\ang{10.79}$, kemiringan pada pertanyaan 2.

\textbf{6.} Serpih tidak menghasilkan gelombang kepala: ia tak terlihat, dan survei
menghitungnya bersama kapur. Di sini $v_1 = \qty{2.40}{km.s^{-1}}$ diberikan sebagai
rata-rata keduanya; seandainya kecepatan gelombang langsung kapur, \qty{2.50}{km.s^{-1}},
dipakai untuk seluruh lapisan penutup, setiap kedalaman akan keluar terlalu besar.

\textbf{7.} $V = 600/\sqrt{0.5437^2 - 0.4878^2} = 600/0.2401 =
\qty{2499}{m.s^{-1}}$ dan $V = 600/\sqrt{0.6240^2 - 0.5748^2} =
\qty{2471}{m.s^{-1}}$. Nilai kedua merata-ratakan kapur dengan serpih yang lebih lambat
di bawahnya: rata-rata turun.

\textbf{8.} \qty{2811}{m.s^{-1}} dan \qty{2829}{m.s^{-1}}.

\textbf{9.} Lapisan 1: \qty{2499}{m.s^{-1}}. Lapisan 2: $v^2 = (2471^2\times0.5748
- 2499^2\times0.4878)/0.0870$, $v = \qty{2308}{m.s^{-1}}$. Lapisan 3:
$v^2 = (2811^2\times0.6548 - 2471^2\times0.5748)/0.0800$,
$v = \qty{4561}{m.s^{-1}}$. Lapisan 4: $v^2 = (2829^2\times0.7348 -
2811^2\times0.6548)/0.0800$, $v = \qty{2972}{m.s^{-1}}$.

\textbf{10.} Kapur (\qty{2499}{m.s^{-1}}), serpih (\qty{2308}{m.s^{-1}}),
batu gamping (\qty{4561}{m.s^{-1}}), batu pasir (\qty{2972}{m.s^{-1}}), masing-masing berada dalam
kisaran yang dipublikasikan. Refleksi-refleksinya adalah antarmuka kapur--serpih, puncak
batu gamping, antarmuka batu gamping--batu pasir, dan ketidakselarasan pada
puncak basemen sekis.

\textbf{11.} Kapur dan serpih memiliki impedansi hampir sama: $R = 0.028$
(latihan 1), dibandingkan 0.38 untuk puncak batu gamping.

\textbf{12.} Hiperbola kecepatan $V_{\text{rms}}$ hanya berlaku hingga orde kedua
dalam offset; pada offset \qty{2}{km}, dua kali kedalaman reflektor terdalam,
waktu-waktunya menyimpang darinya, dan hiperbola yang dicocokkan dengan seluruh bentangan
akan memberikan kecepatan yang keliru, dan interval Dix bahkan lebih buruk. Offset jauh
juga meregangkan wavelet ketika dikoreksi.

\textbf{13.} Waktu bolak-balik di dalam lapisan: \qty{0.4878}{s}, \qty{0.0870}{s},
\qty{0.0800}{s}, \qty{0.0800}{s}. Ketebalan $v\,\Delta t/2$: \qty{609.5}{m},
\qty{100.4}{m}, \qty{182.4}{m}, \qty{118.9}{m}.

\textbf{14.} Kedalaman di bawah permukaan tanah: \qty{609.5}{m}, \qty{709.9}{m},
\qty{892.3}{m}, \qty{1011.2}{m}. Serpih, batu gamping, dan batu pasir keluar
dalam \qty{0.4}{m}, \qty{2.4}{m}, dan \qty{1.1}{m} dari ketebalannya, galat
yang membesar pada batu gamping yang cepat dan tipis, tempat galat waktu kecil
paling merugikan.

\textbf{15.} $\bar v = 2\times1011.2/0.7348 = \qty{2752}{m.s^{-1}}$. Kecepatan RMS
memberi bobot tiap lapisan dengan $v^2$ dan tidak pernah lebih kecil daripada rata-rata
berbobot waktu: $2829\times0.7348/2 = \qty{1039}{m}$, \qty{28}{m} terlalu dalam.

\textbf{16.} $\lambda = 2308/30 = \qty{76.9}{m}$, $\lambda/4 = \qty{19.2}{m}$;
taksiran Ricker adalah $2308/(5.2\times30) = \qty{14.8}{m}$. Serpih
setebal \qty{100}{m}, lima kali \omterm{def:b3:exploration-seismology:resolution}{ketebalan tuning-nya}: puncak dan dasarnya
dapat dibedakan. Ia akan mengalami tuning, menurut aturan seperempat panjang gelombang, pada $f = v/(4b) =
2308/400 = \qty{5.77}{Hz}$.

\textbf{17.} $\lambda = 2829/30 = \qty{94.3}{m}$, $z = \qty{1011}{m}$:
$r_F \approx \sqrt{94.3\times1011/2} = \qty{218}{m}$ (\qty{220}{m} dengan
suku $\lambda^2/16$).

\textbf{18.} $t(x)^2 = t_0^2 + (2x/v)^2$, sehingga $v = 2x/\sqrt{t^2 - t_0^2} =
400/\sqrt{0.5732^2 - 0.550^2} = 400/0.1614 = \qty{2478}{m.s^{-1}}$.

\textbf{19.} $z = vt_0/2 = 2478\times0.550/2 = \qty{681}{m}$.

\textbf{20.} Kedalaman yang diplot $vt/2$ berubah sebesar $1.24\times0.167 = 0.207$ per
satuan panjang: $\theta_a = \arctan0.207 = \ang{11.70}$. Kemiringan sebenarnya memiliki
$\sin\theta = 0.207$: $\theta = \ang{11.95}$, kemiringan sayap.

\textbf{21.} $d = 2480\times0.655/2 = \qty{812}{m}$; migrasi memindahkan
refleksi $d\sin\theta = 812\times0.207 = \qty{168}{m}$ ke atas kemiringan, menuju
puncak antiklin, dan naik dari \qty{812}{m} menjadi $d\cos\theta =
\qty{795}{m}$.

\textbf{22.} $2500\times0.160/2 = \qty{200}{m}$.

\textbf{23.} Dari pertanyaan 13 dan 14, ketidakselarasan basemen terletak
\textbf{\qty[text-series-to-math=true]{1011}{m}} di bawah permukaan tanah pada CMP S.
\end{solution}
